\documentclass[border=6pt]{standalone}
\usepackage[UTF8]{ctex}
\usepackage{tikz}
\usetikzlibrary{positioning, arrows.meta, fit, backgrounds, calc}

\definecolor{mfill}{HTML}{DAEAF6}
\definecolor{medge}{HTML}{2E75B6}
\definecolor{dfill}{HTML}{FFF2CC}
\definecolor{dedge}{HTML}{BF8F00}
\definecolor{pfill}{HTML}{E2EFDA}
\definecolor{pedge}{HTML}{548235}
\definecolor{ffill}{HTML}{F8D7DA}
\definecolor{fedge}{HTML}{C0504D}
\definecolor{synthbg}{HTML}{F2F2F2}
\definecolor{ac}{HTML}{404040}

\begin{document}
\begin{tikzpicture}[
    >=Stealth,
    data/.style={draw=dedge, fill=dfill, rounded corners=2pt,
        minimum height=0.9cm, align=center, font=\small, inner sep=4pt},
    model/.style={draw=medge, fill=mfill, rounded corners=4pt, line width=1pt,
        minimum height=1.3cm, align=center, font=\small, inner sep=5pt},
    proc/.style={draw=pedge, fill=pfill, rounded corners=2pt,
        minimum height=0.9cm, align=center, font=\small, inner sep=4pt},
    filt/.style={draw=fedge, fill=ffill, rounded corners=3pt,
        align=center, font=\small, inner sep=5pt},
    tbox/.style={draw=black!30, fill=white, rounded corners=2pt,
        align=center, font=\small, inner sep=5pt},
    arr/.style={->, line width=0.8pt, ac},
    darr/.style={->, line width=0.8pt, ac, densely dashed},
    slabel/.style={font=\normalsize\bfseries},
    % uniform width for synthesis pipeline nodes
    sp/.style={proc, minimum width=2.2cm, minimum height=1.3cm},
    sd/.style={data, minimum width=2.2cm, minimum height=1.3cm},
]

% ============================================================
% BOTTOM: Inference Pipeline
% ============================================================
\node[data] (inp) at (0,0) {时序图\\图像};
\node[model, right=0.6cm of inp] (m1) {第一阶段\\{\footnotesize Qwen3-VL-4B-Instruct}\\{\footnotesize (多模态大语言模型)}};
\node[data, right=0.6cm of m1] (wjm) {WaveJSON\\结构化表示};
\node[model, right=0.6cm of wjm] (m2) {第二阶段\\{\footnotesize Qwen3-8B}\\{\footnotesize (文本大语言模型)}};
\node[data, right=0.6cm of m2] (outp) {Verilog\\代码};

\draw[arr] (inp) -- (m1);
\draw[arr] (m1) -- (wjm);
\draw[arr] (wjm) -- (m2);
\draw[arr] (m2) -- (outp);

% Geometric center of the whole figure
\path let \p1=(inp.west), \p2=(outp.east) in
    coordinate (fig_cx) at ({(\x1+\x2)/2}, 0);

\node[slabel] at ($(fig_cx)+(0,-1.1)$) {推理阶段};

% ============================================================
% LAYER 2: Training Data Pairs
% ============================================================
\node[tbox, above=1.3cm of m1] (t1) {SFT训练\\(图像, WaveJSON) 对};
\node[tbox, above=1.3cm of m2] (t2) {SFT + RLVR训练\\(WaveJSON, Verilog) 对};

\draw[darr] (t1) -- (m1);
\draw[darr] (t2) -- (m2);

% ============================================================
% LAYER 3: Quality Filter — centered at fig_cx
% ============================================================
\coordinate (filt_y_ref) at ($(t1.north)!0.5!(t2.north)+(0,1.0)$);
\node[filt] (filt) at (fig_cx |- filt_y_ref) {数据质量过滤 + 防泄漏数据划分};

\draw[arr] (filt.south) -- ++(0,-0.35) -| (t1.north);
\draw[arr] (filt.south) -- ++(0,-0.35) -| (t2.north);

% ============================================================
% LAYER 4: Triplet annotation — centered at fig_cx
% ============================================================
\node[font=\small, align=center] (triplet)
    at ($(fig_cx |- filt.north)+(0,0.7)$)
    {训练数据三元组：（时序图图像，WaveJSON，Verilog源代码）};
\draw[arr] (triplet.south) -- (filt.north);

% ============================================================
% TOP: Data Synthesis Pipeline
% All nodes use uniform minimum width (2.2cm) so p3 is the
% true geometric center of the pipeline.
% ============================================================
\coordinate (synth_ref) at ($(fig_cx |- triplet.north)+(0,1.6)$);

\node[sp] (p3) at (synth_ref) {格式转换\\{\tiny(vcd2wavedrom)}};
\node[sp, left=0.25cm of p3]  (p2) {功能仿真\\{\tiny(Icarus Verilog)}};
\node[sp, left=0.25cm of p2]  (p1) {测试平台\\生成};
\node[sd, left=0.25cm of p1]  (vsrc) {Verilog\\源代码};
\node[sd, right=0.25cm of p3] (wj0) {WaveJSON};
\node[sp, right=0.25cm of wj0] (p4) {图像渲染\\{\tiny(wavedrom-cli)}};
\node[sd, right=0.25cm of p4] (im0) {时序图\\图像};

\draw[arr] (vsrc) -- (p1);
\draw[arr] (p1) -- (p2);
\draw[arr] (p2) -- (p3);
\draw[arr] (p3) -- (wj0);
\draw[arr] (wj0) -- (p4);
\draw[arr] (p4) -- (im0);

\begin{scope}[on background layer]
    \node[fill=synthbg, rounded corners=6pt, inner sep=10pt,
        fit=(vsrc)(im0)] (sbox) {};
\end{scope}
\node[slabel, anchor=south] at (sbox.north) {数据合成链路};

% Vertical arrow at fig_cx — guaranteed straight down
\draw[arr] (sbox.south -| fig_cx) -- (triplet.north);

\end{tikzpicture}
\end{document}
